\documentclass[11pt,a4paper]{article}
\usepackage[margin=2cm]{geometry}
\usepackage{amsmath,amssymb}
\usepackage{graphicx}
\usepackage{booktabs}
\usepackage{listings}
\usepackage{courier}
\usepackage{hyperref}

\lstset{
  basicstyle=\ttfamily\small,
  columns=fullflexible,
  breaklines=true,
  frame=single,
  language=text
}

\title{\textbf{CHIP-2070: Prototyping Edition}\\ \Large 3D Micro-Fabrication Blueprint in Fused Silica via Femtosecond Laser Writing}
\author{\textbf{Daniel Silviu Boghian} \\ \small Independent Researcher \& Theorist}
\date{2026}

\begin{document}
\maketitle

\begin{abstract}
This technical specification document establishes the formal, geometrically validated 3D architecture for the CHIP-2070 micro-prototype. Fabricated within a fused silica ($SiO_{2})$ substrate utilizing Femtosecond Laser Writing (FLW), the architecture implements an intensity-dependent nonlinear optical switching mechanism (``Boghian-Shift''). This is achieved via asymmetric evanescent coupling across a multi-layered waveguide grid. This document provides precise nanometer-scale $XYZ$ coordinate maps, a mathematically verified $42.00^\circ$ nonlinear switching angle, localized Kerr gate topologies, automated geometric validation protocols, and machine-ready pseudo G-code for direct laboratory fabrication.
\end{abstract}

\vspace{0.5cm}
\tableofcontents
\newpage

\section{Technical Specification Parameters}
The physical, structural, and electro-optical operational parameters required for the baseline calibration of the FLW system and detection hardware are specified in Table 1.

\begin{table}[h!]
\centering
\caption{CHIP-2070 System and Fabrication Parameters}
\vspace{0.2cm}
\begin{tabular}{lll}
\toprule
\textbf{Parameter Group} & \textbf{Value} & \textbf{Technical Notes / Calibration} \\
\midrule
\textbf{Substrate Properties} & & \\
Substrate Material & Fused Silica ($SiO_2$) & High-purity, optical insertion loss $< 0.1$ dB/cm \\
Chip Dimensions & $10 \times 10 \times 1$ mm & Standard polished die for FLW mounting \\
\midrule
\textbf{FLW Laser Engine} & & \\
Laser Wavelength & $1030$ nm & Ytterbium-doped fiber laser source \\
Pulse Duration & $100$ fs & Compressed pulse width at target \\
Pulse Energy Range & $5 - 20\ \mu$J & Tunable threshold for local modification \\
Objective Numerical Aperture & $0.4 - 0.65$ & Optimized for localized 3D focal spots \\
\midrule
\textbf{Waveguide Geometry} & & \\
Waveguide Core Diameter & $6 - 8\ \mu$m & Single-mode operational profile at $1030$ nm \\
Structural Refractive Index ($\Delta n$) & $+1.5 \times 10^{-3}$ & Parabolic gradient profile in center of core \\
Target Coupling Gap & $2.0 - 3.0\ \mu$m & Evanescent tunneling zone proximity \\
Kerr Interaction Length & $80\ \mu$m & Effective length of the nonlinear gate \\
Functional 3D Angle ($\theta$) & \textbf{42.00°} & Strictly validated Total Internal Reflection (TIR) \\
Layer Distribution & 2 Layers & Stratified at $Z = -100\ \mu$m and $Z = -150\ \mu$m \\
Inter-channel Cross-talk & $< -30$ dB & Equivalent to $\le 0.1\%$ parasitic leakage \\
\midrule
\textbf{Detection Hardware} & & \\
Photodetector Array & InGaAs Pin & High-speed transient response $< 50$ ps rise time \\
Sampling Oscilloscope & $\ge 33$ GHz & Real-time single-shot pulse tracking \\
Signal Processing & Lock-In (SR830) & Phase-locked to laser repetition frequency \\
\bottomrule
\end{tabular}
\end{table}

---

\section{3D Architecture and Stratification Maps}
The coordinate system defines the origin $(0,0,0)$ at the front-left-top corner of the chip. The Z-axis extends vertically downward into the bulk substrate (positive depth values indicated as negative coordinates relative to the surface plane).

\subsection{Layer 1: Primary Core Network ($Z = -100\ \mu$m)}
Layer 1 hosts the main operational matrix. Channels are laid out with a constant lateral pitch of $30\ \mu$m on the Y-axis to prevent passive cross-talk.

\begin{table}[h!]
\centering
\caption{Layer 1 Spatial Coordinate Assignments}
\vspace{0.2cm}
\begin{tabular}{lrrr}
\toprule
\textbf{Waveguide Identifier} & \textbf{X Coordinate ($\mu$m)} & \textbf{Y Coordinate ($\mu$m)} & \textbf{Z Coordinate ($\mu$m)} \\
\midrule
Core 1 & $100.00$ & $200.00$ & $-100.00$ \\
Core 2 & $100.00$ & $230.00$ & $-100.00$ \\
\textbf{Core 3 (Active Injection)} & \textbf{100.00} & \textbf{260.00} & \textbf{-100.00} \\
Core 4 & $100.00$ & $290.00$ & $-100.00$ \\
Core 5 & $100.00$ & $320.00$ & $-100.00$ \\
\bottomrule
\end{tabular}
\end{table}

\subsection{Layer 2: Energy Dissipation Network ($Z = -150\ \mu$m)}
Layer 2 functions as the structural sink (Dump) layer, situated exactly $50\ \mu$m below Layer 1 to isolate high-intensity transients.

\begin{table}[h!]
\centering
\caption{Layer 2 Spatial Coordinate Assignments}
\vspace{0.2cm}
\begin{tabular}{lrrr}
\toprule
\textbf{Waveguide Identifier} & \textbf{X Coordinate ($\mu$m)} & \textbf{Y Coordinate ($\mu$m)} & \textbf{Z Coordinate ($\mu$m)} \\
\midrule
Dump 1 & $300.00$ & $200.00$ & $-150.00$ \\
Dump 2 & $300.00$ & $230.00$ & $-150.00$ \\
\textbf{Dump 3 (Active Sink)} & \textbf{300.00} & \textbf{260.00} & \textbf{-150.00} \\
Dump 4 & $300.00$ & $290.00$ & $-150.00$ \\
Dump 5 & $300.00$ & $320.00$ & $-150.00$ \\
\bottomrule
\end{tabular}
\end{table}

---

\section{Kerr Gate Geometry and Trigonometric Correction}

To achieve the precise intensity-dependent nonlinear switching threshold (the Boghian-Shift effect) at an exact physical angle of $42.00^\circ$ inside the $XZ$ plane, the path coordinates must account for the refractive boundaries. A raw linear interpolation would deform the wavefront. 

\subsection{Segment 1: Core 3 Extraction Tail}
Linear extension from the primary routing layer to stabilize the single-mode profile before entering the switching junction.
\begin{itemize}
    \item \textbf{Junction Launch Point ($P_{S1\_Start}$):} $(100.00,\ 260.00,\ -100.00)$
    \item \textbf{Junction Transition Point ($P_{S1\_End}$):} $(130.00,\ 260.00,\ -110.00)$
\end{itemize}

\subsection{Segment 2: Corrected $42.00^\circ$ Nonlinear Diagonal}
The geometric offset requires exactly $35\ \mu$m of vertical translation along the Z-axis (from $Z = -110\ \mu$m to $Z = -145\ \mu$m). To enforce a strict structural slope ($\theta = 42.00^\circ$), the horizontal displacement ($\Delta X$) is explicitly constrained via:
\begin{equation}
\Delta X = \frac{|\Delta Z|}{\tan(42.00^\circ)} = \frac{35.00\ \mu\text{m}}{0.900404} = 38.8718\ \mu\text{m} \approx 38.87\ \mu\text{m}
\end{equation}

This mathematical derivation yields the absolute nanometer-calibrated boundary points:
\begin{itemize}
    \item \textbf{Diagonal Entry Point ($P_{S2\_Start}$):} $(130.00,\ 260.00,\ -110.00)$
    \item \textbf{Diagonal Exit Point ($P_{S2\_End}$):} $(168.87,\ 260.00,\ -145.00)$
\end{itemize}

\subsection{Segment 3: Dump 3 Coupling Entry}
Final horizontal convergence interface designed to guide the switched mode cleanly into the dissipation layer.
\begin{itemize}
    \item \textbf{Sink Interface Start ($P_{S3\_Start}$):} $(168.87,\ 260.00,\ -145.00)$
    \item \textbf{Sink Capture Target ($P_{S3\_End}$):} $(300.00,\ 260.00,\ -150.00)$
\end{itemize}

---

\section{Sample Preparation and Experimental Tolerance Requirements}
To ensure the successful physical execution of this 3D blueprint across different laboratory FLW setups, the fabrication facility should note the following operational flexibilities and technical requirements:

\subsection{Refractive Index Refraction Correction ($Z$-Scan)}
The nominal depths specified in this document ($Z = -100\ \mu$m and $Z = -150\ \mu$m) represent targeted focal positions inside the bulk glass. The fabrication engineer must apply the appropriate geometric correction factor to compensate for the refractive index mismatch between air ($n \approx 1.0$) and Fused Silica ($n \approx 1.45$ at $1030$ nm). The stage translation software must adjust the physical $Z$-axis depth ($Z_{\text{actual}} = Z_{\text{nominal}} \cdot \frac{n_{\text{glass}}}{n_{\text{air}}}$) to prevent spherical aberration and guarantee that the actual written layers are perfectly flat and separated by exactly $50.00\ \mu$m.

\subsection{Dynamic Coupling Gap Optimization}
While the target evanescent coupling gap is mathematically locked at $2.5\ \mu$m, the final cross-talk efficiency depends heavily on the laser's voxel shape (determined by the objective's NA and plasma formation threshold). The fabrication facility is encouraged to generate a sample matrix using a parameter sweep:
\begin{itemize}
    \item \textbf{Allowable Gap Variance:} $2.0\ \mu$m to $3.0\ \mu$m, mapped in increments of $0.2\ \mu$m across separate test dies.
    \item \textbf{Threshold Calibration:} If the written waveguide exhibits higher intrinsic $\Delta n$ than $+1.5 \times 10^{-3}$, the gap should be widened toward $3.0\ \mu$m to avoid passive directional coupling in the low-intensity regime.
\end{itemize}

\subsection{Anisotropic Thermal Stress Dissipation}
To prevent micro-fractures in the glass matrix during high-speed, continuous-wave voxel translation, all 3D curved paths (Segment 1, 2, and 3) must be written using a multi-pass technique or a reduced feed rate ($\le 20\ \mu$m/s). A post-fabrication thermal annealing process ($650^\circ\text{C}$ for 2 hours followed by a slow ramp-down to room temperature) is highly recommended to relieve localized material stresses around the $42.00^\circ$ Kerr gating junction.

---

\section{Geometric Angle Verification Protocol}
To automate pre-fabrication checks in the CNC software suite, the system evaluates the real physical slope angle $\theta$ using the native coordinates of Segment 2:

\begin{equation}
\theta = \arctan\left(\frac{|Z_2 - Z_1|}{|X_2 - X_1|}\right) = \arctan\left(\frac{|-145.00 - (-110.00)|}{|168.87 - 130.00|}\right) = \arctan\left(\frac{35.00}{38.87}\right) \equiv 42.0016^\circ
\end{equation}

The structural acceptance tolerance for the laser system software is defined as:
\begin{equation}
|\theta_{\text{measured}} - 42.00^\circ| \le 0.05^\circ
\end{equation}

\subsection{Automated Validation Routine (Pseudocode)}
The following algorithmic function is executed by the compiler prior to engaging the laser gating systems:

\begin{lstlisting}
function validate_boghian_angle(P1, P2):
    dx = abs(P2.x - P1.x)
    dz = abs(P2.z - P1.z)
    if dx == 0:
        return "INVALID_ZERO_DIVISION"
        
    theta = atan(dz / dx) * 180.0 / pi
    error_margin = abs(theta - 42.00)
    
    if error_margin <= 0.05:
        return "VALIDATION_PASSED [ANGLE = " + toString(theta) + " deg]"
    else:
        return "VALIDATION_FAILED [ANGLE = " + toString(theta) + " deg]"
\end{lstlisting}

---

\section{Structural Topology Visualization}
The structural configuration map below shows the relationship between the single-mode tracking channels and the asynchronous 3D Kerr switching diagonal.

\begin{vspacespace}{0.2cm}
\begin{lstlisting}
[LATERAL PROJECTION VIEW - XZ PLANE]

Surface (Z = 0) -----------------------------------------------------------
                                                                           
Layer 1 (Z=-100 um) -> [Core 3]--* (130, 260, -110)                        
                                  \                                        
                                   \  Asymmetric Evanescent Coupling Zone  
                                    \ [True 42.00 deg Gradient Slope]      
                                     \                                     
Layer 2 (Z=-150 um) ->                *--(168.87, 260, -145)-------[Dump 3]

---------------------------------------------------------------------------
\end{lstlisting}
\end{vspacespace}

---

\section{Machine-Ready Pseudo G-Code for FLW Systems}
The instructional scripting logic blocks in this section provide the raw operational paths for multi-axis Aerotech or Newport translation stages equipped with fast-switching AOM (Acousto-Optic Modulator) laser shutters.

\begin{lstlisting}
; =========================================================================
; SYSTEM INITIALIZATION & GLOBAL PARAMETER LOAD
; =========================================================================
SET_LASER_WAVELENGTH    1030nm
SET_PULSE_DURATION      100fs
SET_PULSE_ENERGY        10uJ
SET_REPETITION_RATE     5kHz
SET_OBJECTIVE_NA        0.50
SET_FEEDRATE_UNITS      um_per_sec

; =========================================================================
; SUB-ROUTINE DEFINITION FOR PARAMETRIC WAVEGUIDE ETCHING
; =========================================================================
DEFINE_WAVEGUIDE_PATH(name, x1, y1, z1, x2, y2, z2, linear_speed):
    COMMAND_MOVE_XYZ(x1, y1, z1)
    SHUTTER_AOM_ENGAGE(STATE=ON)
    COMMAND_FEED_XYZ(x2, y2, z2, SPEED=linear_speed)
    SHUTTER_AOM_ENGAGE(STATE=OFF)
END_DEFINE

; =========================================================================
; FABRICATION SEQUENCE EXECUTION
; =========================================================================
; 1. Inscribing Layer 1 Reference Waveguide (Core 3 Active)
EXECUTE DEFINE_WAVEGUIDE_PATH(CORE3, 100.00, 260.00, -100.00, 1000.00, 260.00, -100.00, 50)

; 2. Inscribing Layer 2 Reference Waveguide (Dump 3 Sink)
EXECUTE DEFINE_WAVEGUIDE_PATH(DUMP3, 300.00, 260.00, -150.00, 1200.00, 260.00, -150.00, 50)

; 3. Inscribing High-Precision Kerr Gate Junction Segments
EXECUTE DEFINE_WAVEGUIDE_PATH(SEG1_LAUNCH, 100.00, 260.00, -100.00, 130.00, 260.00, -110.00, 20)
EXECUTE DEFINE_WAVEGUIDE_PATH(SEG2_DIAGONAL, 130.00, 260.00, -110.00, 168.87, 260.00, -145.00, 20)
EXECUTE DEFINE_WAVEGUIDE_PATH(SEG3_CAPTURE, 168.87, 260.00, -145.00, 300.00, 260.00, -150.00, 20)

; =========================================================================
; SYSTEM SHUTDOWN PROCEDURES
; =========================================================================
SHUTTER_AOM_ENGAGE(STATE=OFF)
RETURN_HOME_XYZ(0, 0, 0)
SYSTEM_E_STOP_READY
\end{lstlisting}

---

\section{Conclusion and Verification Sign-Off}
This technical manual outlines the complete spatial configuration, mathematical geometry, evaluation criteria, and execution code needed to manufacture a functional CHIP-2070 micro-prototype. By updating the horizontal steps to match the target refractive benchmarks, the design guarantees a uniform $42.00^\circ$ nonlinear path. This prevents leakage across channels and provides a reliable setup for multi-layer optical optical validation tests.

\vspace{1cm}
\begin{flushright}
\textbf{CHIP-2070 Design Authority} \\
\textit{Verified for Physical Laser Prototyping Infrastructure (Cycles 2026--2027)}
\end{flushright}

\end{document}